\documentclass[10pt,a4paper]{amsart}
\usepackage[margin=19mm]{geometry}
\usepackage{amsmath,amssymb,mathtools}
\usepackage[scr=rsfs,cal=euler]{mathalfa}
\usepackage{xfrac}
\usepackage[unicode,hidelinks,bookmarks=false]{hyperref}
\newcommand{\ie}{i.e.\ }
\newcommand{\cf}{cf.\ }
\newcommand{\resp}{respectively\ }
\newcommand{\Subsection}{\subsection}
\providecommand{\nobreakdash}{\nobreak\hbox{-}}
\setlength{\emergencystretch}{2em}
\multlinegap0pt
\usepackage{amssymb}
\usepackage{xcolor}
\usepackage[shortlabels]{enumitem}
\newtheorem{definition}{Definition}
\newtheorem{lemma}{Lemma}
\newcommand{\F}{\mathcal{F}(M,P)}
\newcommand{\bbN}{\mathbb{N}}
\newcommand{\bbR}{\mathbb{R}}
\newcommand{\textr}[1]{\textcolor{black}{#1}}
\title{Homotopy type of stabilizers of smooth functions with non-isolated singularities on surfaces}
\author{Bohdan Feshchenko}
\date{Source: arXiv:2305.08255v3 (CC BY 4.0)}
\begin{document}
\maketitle

\noindent\emph{Source and licence.} Homotopy type of stabilizers of smooth functions with non-isolated singularities on surfaces. Version arXiv:2305.08255v3 (CC BY 4.0), distributed by arXiv under the Creative Commons Attribution 4.0 International licence, http://creativecommons.org/licenses/by/4.0/, as recorded in the arXiv licence metadata for this exact version and shown on the abstract page https://arxiv.org/abs/2305.08255v3. Full licence notice: \texttt{licenses/arxiv-2305.08255-CC-BY-4.0.txt}.\par\smallskip

\noindent\textit{Editorial scope.} Counted unit: the article's lemma lm:vector-field-FC (source0.tex lines 668--679). The article's complete original proof of it is delivered with it (source0.tex lines 680--743, one \texttt{proof} environment of 64 lines). No mathematics is written, repaired or re-derived: the statement and the proof are the article's own lines, in the article's own notation, and no retained line is substituted or re-flowed.  Retained uncounted as the source-local context the proof needs: lines 118--129; lines 444--454; lines 586--586; lines 590--593.  Nothing was omitted from any retained span, so the package carries no omission record.  No retained line contains a citation, so the wrapper carries no bibliography and no bibliography entry.  No retained line contains a figure environment, an image-inclusion command or drawing code, so no image is delivered and the package declares no asset dependency.  Editorial formatting only, in addition: an AMS \texttt{amsart} preamble that restates the article's own declarations and macros used by the retained lines, namely the environments \texttt{definition}, \texttt{lemma} on separate counters, together with the standard packages those lines need.  The retained environments print in this document as Definition 1, Lemma 1 in order of appearance, whereas the article prints the same items with the numbering of its own full document.  The complete native source of the article as distributed in the arXiv e-print for this exact version is bundled unchanged as \texttt{sources/141478/source0.tex}.
\begin{definition}\label{def:class-F} 
	 \normalfont A smooth function $f\in C_{\partial}^{\infty}(M,P)$ on $M$ belongs to the class $\mathcal{F}(M, P)$ if it satisfies the following conditions:
	\begin{enumerate}
		\item the set of critical points $\Sigma_f$ of $f$ is a disjoint union of smooth submanifolds of $M$ and $\Sigma_f\subset \mathrm{Int}(M),$
		\item  for each connected component $C$ of $\Sigma_f$ and any critical point $p\in C$,
		the germ $(f,p)$ of $f$ at $p$ is smoothly equivalent 
		\begin{enumerate}
			\item[(a)] to either \textr{the germ at $0\in \bbR^2$ of a homogeneous polynomial} $f_p:\bbR^2\to \bbR$ without multiple factors  with $\deg f_p\geq 2$,
			\item[(b)] or to \textr{the germ at $0\in \bbR^2$ } of  $f_C(x,y) = \pm y^{n_C}$  for some $n_C\in \bbN_{\geq 2}$ depending of $C.$  
		\end{enumerate}
	\end{enumerate}
\end{definition}

\begin{lemma}\label{lm:technical-lema}
	Let $h:\mathbb{R}^2\to \bbR^2$ be a diffeomorphism given by $h(x,y) = (u(x,y), v(y))$ with $u_x >0$, $v'>0$ and $v(0) \equiv 0$. Denote by  $G = h^*(F_{(u,v)})$  a pull-back of $F_{(u,v)}$ along $h$. Then 
	\begin{enumerate}
		\item $G = \tau F_{(x,y)}$ for some positive smooth function $\tau:\bbR^2\to \bbR$,
		\item for any smooth  function $\rho:\bbR^2\to \bbR$ with $0\leq \rho\leq 1$ we have 
		$$
		\rho F_{(x,y)}+(1-\rho) G = \eta F_{(x,y)}
		$$
		 for some smooth and positive function $\eta:\mathbb{R}^2\to \bbR$.
	\end{enumerate}
\end{lemma}

\subsection{Hamiltonian vector fields for functions from $\F$}\label{ssec:ham-function-F}

\begin{align}
	X_f &= \frac{1}{g(x,y)}  (f_z)'_y\frac{\partial}{\partial x} - \frac{1}{g(x,y)} (f_z)'_x\frac{\partial}{\partial y}, &(\text{if $z$ is isolated})\label{eq:ham-vf-near-iso}   \\
	 X_f &= \pm\frac{n_Cy^{n_C-1}}{{g}(x,y)}\frac{\partial}{\partial x},& (\text{if $z$ is not isolated, $z\in C$})\label{eq:Ham-near-C}
\end{align}

\begin{lemma}\label{lm:vector-field-FC}
	Let $M$ be a smooth, compact, connected and \textr{orientable} surface, let $f$ be a function from $\F$ and $C$ be a critical circle of $f$. Then there exists a saturated neighborhood $Q_C$ of $C$ and a vector field $F_C$ on $\overline{Q_C}$ such that 
	\begin{enumerate}
		\item $f$ is constant along $F_C$ on $\overline{Q_C}$,
		\item for any $z\in C$ there exists a local chart $(U,(x,y))$ near $z$ such that $F_C$ on $U$ has the form
		\begin{align}
			F_C &= \frac{\partial}{\partial x},&& \text{if $C$ is non-extremal},\label{eq:FC-non-extremal}\\
			F_C &= \mu_z(x,y)y\frac{\partial }{\partial x}, && \text{if $C$ is extremal}\label{eq:FC-extremal},
		\end{align}
		for some smooth and positive function $\mu_z:U\to \bbR$.
	\end{enumerate}
\end{lemma}

\begin{proof}
Let also $z\in C$ be a non-isolated critical point of $f$. By (2.b) of Definition \ref{def:class-F} there exist a local coordinate system $(U_z, (\tilde{x},\tilde{y}))$ near $z$ such that a local representation of $f$ on $U_z$ is $f_z(\tilde{x},\tilde{y}) = a \tilde{y}^{n_C}$ for  some $n_C\geq 2$ depending on $C$, where $a \in \{ \pm 1\}$.

Let $\omega$ be a symplectic form on $M$.
 A Hamiltonian vector field $X_{f_z}$ of $f_z$ on $U_z$ has the form \eqref{eq:Ham-near-C}. By changing coordinate system $(x,y)= (a\tilde{x},\tilde{y})$ on $U_z$, we obtain that $X_f$ on $U_z$ has the form:
\begin{equation}\label{eq:Ham-chang-coord}
	X_f = \frac{n_Cy^{n_C-1}}{{g}(a x,y)}\frac{\partial}{\partial x},
\end{equation}
where $g:U_z\to \bbR$ is some positive and smooth function, see \textsection\ref{ssec:ham-function-F}.

The collection of sets $\{U_z\,|\,z\in C \}$ forms an open cover of $C$. From compactness of $C$, we can choose a finite ``chain-like'' subcover of $\{U_z\}_{z\in C}$. Namely,  
there exist $N\in \bbN$ and a subset $\{z_0, z_1,\ldots, z_{N-1}\}\subset C$ such that
\begin{enumerate}
	\item $\{U_{z_i}\}_{i = 0}^{N-1}$ is an open cover of $C$,
	\item $U_{z_i}\cap U_{z_j}\neq \varnothing$ iff $j = i\pm 1\;\mathrm{mod}\, N$.  
\end{enumerate}
{Choose a coordinate system $\phi_i = (x_i, y_i):U_{z_i}\to\mathbb{R}^2$ on $U_{z_i}$ with $\phi_i(C\cap U_{z_i})\subset \{y = 0\}$} such that  $X_f$ on $U_{z_i}$ has the form \eqref{eq:Ham-chang-coord} and
define a vector field $F_{z_i}$ on $U_{z_i}$ by the following formulas
\begin{align}
	F_{z_i} &= \frac{\partial}{\partial x_i}, &(\text{if $C$ is non-extremal, i.e., $n_C$ is odd}), \label{eq:F_z-C-non-extr}\\
	F_{z_i} &= y_i\frac{\partial}{\partial x_i}, &(\text{if $C$ is extremal, i.e., $n_C$ is even}). \label{eq:F_z-C-extr}
\end{align}
Note that the flows of vector fields $F_{z_i}$ and $X_f$ define the same foliation on $U_{z_i}\setminus C$.
If $C$ is non-extremal critical circle of $f$, then $U_z\cap C$ is a regular trajectory for $F_{z_i}$. 
From Eq.~\eqref{eq:Ham-chang-coord} it is easy to see that vector fields $F_{z_i}$  and $X_f$ are codirectional on $U_{z_i}\setminus C$.


Let $Q_C$ be a saturated neighborhood of $C$ such that $\overline{Q_C}\subset \bigcup_{i = 0}^{N-1} U_{z_i}.$ We will use the partition of unity to define a  vector field $F_C$ on $\overline{Q}_C$ mentioned above.
We put $Q_i = Q_C\cap U_{z_i}$ and $V_i = Q_i\setminus (U_{z_{i-1}}\cup U_{z_{i+1}})$, where all indexes are taken modulo $N$. It is easy to see that there exists a family smooth functions $\{ \rho_i:\overline{Q_C}\to \bbR\,|\, \rho_i\geq 0,\ i =0,1,\ldots, N-1 \}$
which satisfies the following conditions:
\begin{itemize}
	\item $\rho_i = 1$ on $\overline{V_i}$,
	\item $\rho_i = 0$ on $\overline{Q}\setminus \overline{Q_{i}}$,
	\item $\rho_i(z) + \rho_{i+1}(z) = 1$ for $z\in \overline{Q_i}\cap \overline{Q_{i+1}}$,
\end{itemize}
where  indexes are taken modulo $N$. Define a vector field $F_C$ on $\overline{Q_C}$ by the formula
\begin{equation}\label{eq:F_C-formula}
	F_C =\begin{cases}
		F_{z_i}, &\text{on }\overline{V_i},\\
		\rho_iF_{z_i}+\rho_{i+1} F_{z_{i+1}}, &\text{on } \overline{Q_i\cap Q_{i+1}}.
	\end{cases}
\end{equation}








Clearly, $F_C$ is codirectional with $X_f$ on $\overline{Q_C}\setminus C$ and $F_C(f|_{\overline{Q_C}}) = 0$. Therefore $f$ is constant along $F_C$ on $\overline{Q_C}$.



If $C$ is a non-extremal critical circle of $f$, then $F_C$ has no zeros in $\overline{Q_C}$. Thus,  for each $z\in C$ there exists a coordinate chart $(U, (x,y))$ near $z$ such that $F_C$ on $U$ has the form \eqref{eq:FC-non-extremal}.


{If $C$ is an extremal critical circle of $f$, then on $\overline{V_i}$ a vector field $F_C$ has the form $F_C = F_{z_i} = y_i\frac{\partial}{\partial x_i}$. It remains to prove that a vector field 
$\rho_iF_{z_i}+\rho_{i+1} F_{z_{i+1}}$ as above has the form $\eta_{i} F_{z_{i}}$ for some positive smooth function $\eta_{i}:\overline{Q_i\cap Q_{i+1}}\to \bbR$. This is a consequence of  Lemma \ref{lm:technical-lema}.

Indeed,  vector fields $F_{z_i}$ and $F_{z_{i+1}}$ have the same trajectories on $U_{z_i}\cap U_{z_{i+1}}$, and they are codirectional on $(U_{z_i}\cap U_{z_{i+1}})\setminus C$. Then there exist a ``special '' coordinate change, i.e., an orientation-preserving diffeomorphism 
$h:\phi_i(\overline{Q_i\cap Q_{i+1}})\to \phi_{i+1}(\overline{Q_i\cap Q_{i+1}})$ which satisfies conditions of the Lemma \ref{lm:technical-lema}, i.e., $h$ maps  trajectories of $(\phi_i)_*F_{z_i}$ to  trajectories of $(\phi_{i+1})_*F_{z_{i+1}}$ and preserve orientations of their regular trajectories. Thus, by (2) of Lemma \ref{lm:technical-lema}, on $\overline{Q_i\cap Q_{i+1}}$ a vector field
$\rho_iF_{z_i}+\rho_{i+1} F_{z_{i+1}}$ has the form $\eta_{i} F_{z_{i}}$ for some positive smooth function $\eta_{i}:\overline{Q_i\cap Q_{i+1}}\to \bbR$. This completes the proof of  Lemma \ref{lm:vector-field-FC}.}
\end{proof}

\end{document}
